\documentclass[11pt]{article}

\usepackage[T1]{fontenc}
\usepackage{lmodern}
\usepackage[margin=1in]{geometry}
\usepackage{amsmath,amssymb,amsthm}
\usepackage{booktabs}
\usepackage{microtype}
\usepackage[hidelinks]{hyperref}
\setlength{\emergencystretch}{2em}

\newtheorem{theorem}{Theorem}[section]
\newtheorem{proposition}[theorem]{Proposition}
\newtheorem{lemma}[theorem]{Lemma}

\newcommand{\Pvec}{\mathcal P}
\newcommand{\N}{\mathbb N}

\hypersetup{
  pdftitle={Sharp Initial Thresholds for Abelian-Bordered Binary Infinite Words},
  pdfauthor={Carptopus},
  pdfsubject={Abelian-bordered factors and ultimate periodicity of binary infinite words},
  pdfkeywords={Abelian border, infinite word, ultimate periodicity, combinatorics on words, finite overlap graph, computer-assisted proof}
}

\title{Sharp Initial Thresholds for Abelian-Bordered\\
Binary Infinite Words}
\author{Carptopus\\\small\texttt{carptopus@163.com}}
\date{26 August 2026\\\small v0.1-beta; not peer reviewed}

\begin{document}
\maketitle

\begin{abstract}
Let $\mu_{\mathrm{ab}}(x)$ denote the maximum length of an
Abelian-unbordered factor of a binary infinite word $x$, with value $\infty$
when these lengths are unbounded. Charlier, Harju, Puzynina and Zamboni asked
whether $\mu_{\mathrm{ab}}(x)<\infty$ forces $x$ to be Abelian ultimately
periodic. We determine the first sharp threshold layers. If
$\mu_{\mathrm{ab}}(x)\le13$, then $x$ is ultimately periodic in the ordinary
sense and its eventual period can be chosen of length at most $13$. Both the
threshold and the period bound are sharp: the least finite value of
$\mu_{\mathrm{ab}}$ attained by a non-ultimately-periodic binary word is $14$.
If $\mu_{\mathrm{ab}}(x)\le14$, then $x$ is Abelian ultimately periodic with
Abelian period length at most $14$, again sharply.

The upper bounds are computer-assisted but exact. We classify all recurrent
cores of two finite overlap graphs. At threshold $14$, every cyclic strongly
connected component is a directed cycle of length at most $13$. At threshold
$15$, all but four cyclic components are directed cycles, and the four
branching components are absorbed by closed Parikh phases of lengths $4$ or
$14$. A separate exact computation establishes the sharp period bounds.

We also prove a non-computational sufficient condition valid at arbitrary
thresholds. If the frequency of $1$ in $x$ is the reduced fraction $p/q$,
$M=\mu_{\mathrm{ab}}(x)<\infty$, and $M<2q$, then $x$ is Abelian ultimately
periodic with blocks of length $q$ and weight $p$. The general question
remains open outside these regimes.
\end{abstract}

\noindent\textbf{Keywords.} Abelian border; infinite word; ultimate
periodicity; combinatorics on words; finite overlap graph; computer-assisted
proof.

\medskip
\noindent\textbf{2020 Mathematics Subject Classification.} 68R15, 05A05.

\section{Introduction}

A nonempty finite word is bordered when it has a nonempty proper prefix equal
to a suffix. Ehrenfeucht and Silberger~\cite{ES} proved that an infinite word
is purely periodic exactly when it has only finitely many unbordered factors.
Replacing equality by Abelian equivalence gives a natural commutative
analogue: two words are Abelian equivalent when they have the same Parikh
vector, and a word is Abelian bordered when a nonempty proper prefix is
Abelian equivalent to a suffix.

Charlier, Harju, Puzynina and Zamboni~\cite{CHPZ} asked whether a binary
infinite word having only finitely many Abelian-unbordered factors must be
Abelian ultimately periodic. They constructed aperiodic words all of whose
factors of length at least $15$ are Abelian bordered, but their examples
remain Abelian periodic. The problem was still listed as open in the survey
of Fici and Puzynina~\cite[Problem~42]{FP}.

For a binary infinite word $x$, define
\[
\mu_{\mathrm{ab}}(x)=\sup\{\,|u|:u\text{ is an Abelian-unbordered factor of }x\,\}.
\]
Thus the open problem asks whether $\mu_{\mathrm{ab}}(x)<\infty$ implies
Abelian ultimate periodicity. Our first result determines the exact initial
boundary at which ordinary ultimate periodicity can fail.

\begin{theorem}\label{thm:ordinary}
If $\mu_{\mathrm{ab}}(x)\le13$, then $x$ is ultimately periodic and its
eventual period can be chosen of length at most $13$. Both constants are
sharp. In particular,
\[
\min\{\mu_{\mathrm{ab}}(x):x\text{ is not ultimately periodic and }
\mu_{\mathrm{ab}}(x)<\infty\}=14.
\]
\end{theorem}

The first layer permitting ordinary aperiodicity still satisfies the desired
Abelian conclusion.

\begin{theorem}\label{thm:abelian}
If $\mu_{\mathrm{ab}}(x)\le14$, then $x$ is Abelian ultimately periodic with
Abelian period length at most $14$. The period bound $14$ is sharp.
\end{theorem}

Theorems~\ref{thm:ordinary} and~\ref{thm:abelian} are proved by finite
certificates which cover every infinite word satisfying the corresponding
threshold condition; they are not bounded-prefix experiments. Our second
type of result comes from the rational frequency and bounded discrepancy
forced by the hypothesis.

\begin{theorem}\label{thm:denominator}
Let $M=\mu_{\mathrm{ab}}(x)<\infty$, and write the frequency of $1$ as the
reduced fraction $p/q$. If $M<2q$, then $x$ is Abelian ultimately periodic.
More precisely, after deleting a finite prefix, $x$ factors into consecutive
blocks of length $q$, each containing exactly $p$ occurrences of $1$.
\end{theorem}

Theorem~\ref{thm:denominator} covers an infinite parameter regime but not the
whole open problem. The remaining difficulty begins when the extremal
discrepancy layer can have several possible return lengths.

\section{Preliminaries}

Let $x=x_0x_1x_2\cdots\in\{0,1\}^{\N}$. For a finite binary word $u$, its
Parikh vector is
\[
\Pvec(u)=(|u|_0,|u|_1).
\]
Words $u,v$ are Abelian equivalent, written $u\sim_{\mathrm{ab}}v$, when
$\Pvec(u)=\Pvec(v)$. A word $u$ of length $L$ has an Abelian border of length
$k$ if $1\le k<L$ and its prefix and suffix of length $k$ are Abelian
equivalent. If $k>L/2$, deleting the common overlap from the two ends produces
an Abelian border of length $L-k<L/2$. It is therefore enough to consider
\[
1\le k\le\lfloor L/2\rfloor.
\]

An infinite word is ultimately periodic if it has the form $uv^\omega$. It
is Abelian ultimately periodic if it admits a factorization
\[
x=uv_0v_1v_2\cdots,
\qquad v_i\sim_{\mathrm{ab}}v_j\quad(i,j\ge0).
\]
The common blocks consequently have the same length and the same number of
ones. We call that common length an Abelian period length.

For integers $C\le N$, call a binary word of length $N$ $(C,N)$-admissible if
every one of its factors whose length lies in $[C,N]$ is Abelian bordered.

\section{Finite overlap certificates}

\subsection{The overlap graph}

Define a finite directed labeled graph $G_{C,N}$ as follows. Its vertices are
all length-$(N-1)$ prefixes and suffixes of $(C,N)$-admissible words. Every
admissible word $w=w_0w_1\cdots w_{N-1}$ gives the edge
\[
w_0\cdots w_{N-2}\longrightarrow w_1\cdots w_{N-1},
\]
labeled by $w_{N-1}$.

If every factor of an infinite word $x$ of length at least $C$ is Abelian
bordered, then its consecutive length-$N$ windows form an infinite path in
$G_{C,N}$. Since the condensation of a finite directed graph is acyclic, the
path eventually remains inside a strongly connected component containing a
directed cycle.

\begin{lemma}\label{lem:cycle}
Let $H$ be a finite strongly connected directed simple graph. If the number
of internal edges of $H$ equals the number of vertices, then $H$ is a directed
cycle.
\end{lemma}

\begin{proof}
Every vertex of a strongly connected graph has positive internal indegree and
outdegree. The sums of the internal indegrees and outdegrees both equal the
number of edges. Equality with the number of vertices forces every internal
indegree and outdegree to be one. Thus $H$ is a directed cycle.
\end{proof}

An infinite path which eventually remains in such a component has an
ultimately periodic label sequence.

\subsection{Closed Parikh phases}

Branching components require a weaker certificate. Let $H$ be a strongly
connected component, let $\ell\ge1$, and let $0\le m\le\ell$. A set
$Q_{\ell,m}\subseteq V(H)$ is a closed $(\ell,m)$-phase if every
length-$\ell$ path in $H$ beginning in $Q_{\ell,m}$ has label weight $m$ and
terminates in $Q_{\ell,m}$.

\begin{lemma}\label{lem:phase}
Every infinite path in $H$ which visits $Q_{\ell,m}$ factors from that visit
onward into label blocks of length $\ell$ and weight $m$.
\end{lemma}

\begin{proof}
Apply the two defining properties successively to consecutive length-$\ell$
path segments. Each segment has the same binary Parikh vector
$(\ell-m,m)$, and its terminal state is again in the phase.
\end{proof}

For fixed $(\ell,m)$, the verification program computes the greatest closed
phase by iterating the monotone operator
\[
\Phi_{\ell,m}(Q)=\{v\in V(H):\text{ every length-}\ell\text{ path from }v
\text{ has weight }m\text{ and ends in }Q\}
\]
downward from $V(H)$. Stabilization is exact because $H$ is finite. Several
closed phases may be used at once.

\begin{lemma}\label{lem:absorb}
Suppose $Q$ is a union of closed Parikh phases in $H$, and the subgraph
induced by $V(H)\setminus Q$ is acyclic. Then every infinite path in $H$ is
Abelian ultimately periodic.
\end{lemma}

\begin{proof}
An infinite path cannot remain forever in a finite acyclic graph, so it
eventually enters one of the closed phases contained in $Q$.
Lemma~\ref{lem:phase} then gives an Abelian periodic factorization of its
label tail.
\end{proof}

These lemmas reduce both upper bounds to exact finite graph statements.

\section{Exact classification of the first two layers}

The admissible words and overlap graphs were constructed independently in
C++ and Python. The Python verification also exhaustively compares the two
definitions of Abelian borderedness on all $131{,}070$ nonempty binary words
of length at most $16$, and checks small overlap graphs by brute force.

\begin{proposition}\label{prop:c14}
For $(C,N)=(14,36)$, the graph $G_{14,36}$ has
\[
\begin{array}{lr}
\text{admissible length-36 words} & 38{,}450,\\
\text{vertices} & 38{,}544,\\
\text{cyclic strongly connected components} & 1{,}377.
\end{array}
\]
Every cyclic strongly connected component is a directed cycle, and the
largest such cycle has $13$ states.
\end{proposition}

\begin{proof}[Proof of the upper bound in Theorem~\ref{thm:ordinary}]
If $\mu_{\mathrm{ab}}(x)\le13$, every factor of length at least $14$ is
Abelian bordered. The length-$36$ windows of $x$ therefore form an infinite
path in $G_{14,36}$. Eventually the path remains in a cyclic strongly
connected component. Proposition~\ref{prop:c14} and Lemma~\ref{lem:cycle}
show that this component is a directed cycle of length at most $13$. Hence
$x$ is ultimately periodic with eventual period length at most $13$.
\end{proof}

\begin{proposition}\label{prop:c15}
For $(C,N)=(15,30)$, the graph $G_{15,30}$ has
\[
\begin{array}{lr}
\text{admissible length-30 words} & 199{,}420,\\
\text{vertices} & 217{,}280,\\
\text{cyclic strongly connected components} & 2{,}540.
\end{array}
\]
Exactly four cyclic components branch. Two have $40$ vertices and $42$
edges; in each, closed phases include a phase of length $4$ and weight $2$,
their union covers $29$ vertices, and the uncovered induced subgraph is
acyclic. The other two have $60$ vertices and $62$ edges; closed phases
include a phase of length $14$ and weight $7$, their union covers $54$
vertices, and the uncovered induced subgraph is acyclic.

The other $2{,}536$ cyclic components are directed cycles. A separate exact
cycle-label post-processing computation shows that every one of these cycles
has an Abelian period length at most $14$.
\end{proposition}

\begin{proof}[Proof of the upper bound in Theorem~\ref{thm:abelian}]
If $\mu_{\mathrm{ab}}(x)\le14$, the length-$30$ windows of $x$ form an
infinite path in $G_{15,30}$. If its eventual component is one of the four
branching components, Proposition~\ref{prop:c15} and Lemma~\ref{lem:absorb}
give a block factorization of length $4$ or $14$. Otherwise the eventual
component is a directed cycle whose label word has an Abelian period of
length at most $14$. In all cases $x$ is Abelian ultimately periodic with an
Abelian period length at most $14$.
\end{proof}

The numerical assertions in Propositions~\ref{prop:c14} and~\ref{prop:c15}
are certificate-bearing finite claims. The verification entry points and
fixed outputs are described in Section~\ref{sec:verification}.

\section{Sharpness}

\subsection{The threshold for ordinary periodicity}

Charlier et al.~\cite[Proposition~2 and Remark~2]{CHPZ} use the two codewords
\[
a=010100110011,
\qquad b=0101001100110011.
\]
Any infinite concatenation of $a$ and $b$ has every factor of length at least
$15$ Abelian bordered. Choosing a non-ultimately-periodic codeword selector
produces a word which is not ultimately periodic: in the aligned length-$4$
block alphabet, the block $0101$ marks the beginning of a codeword and is
followed by two or three copies of $0011$, so the selector can be recovered
from the binary word.

The same construction contains the factor
\[
u=01001100110011
\]
of length $14$. For border lengths $k=1,\ldots,7$, the numbers of ones in the
corresponding prefixes and suffixes are respectively
\[
(0,1),(1,2),(1,2),(1,2),(2,3),(3,4),(3,4).
\]
Thus no such $k$ is an Abelian border. A border longer than $7$ would, after
deleting the common overlap, give a shorter Abelian border. Hence $u$ is
Abelian unbordered and the constructed word has $\mu_{\mathrm{ab}}=14$.
Together with Theorem~\ref{thm:ordinary} this proves the exact threshold
assertion.

\subsection{Sharp period bounds}

Let $w=10^{12}$. The periodic word $w^\omega$ has least ordinary period $13$:
for every $1\le p<13$, its symbol at position $0$ differs from the symbol at
position $p$. Every factor of length $L\ge14$ has the ordinary border of
length $L-13$. On the other hand, $w$ is Abelian unbordered because for
$1\le k\le6$ its prefix of length $k$ has one occurrence of $1$, while its
suffix has none. Therefore
\[
\mu_{\mathrm{ab}}(w^\omega)=13,
\]
and the period-length bound in Theorem~\ref{thm:ordinary} is sharp.

Similarly, put $v=10^{13}$. Then $v$ is Abelian unbordered, while every factor
of $v^\omega$ of length at least $15$ is ordinarily bordered. Hence
\[
\mu_{\mathrm{ab}}(v^\omega)=14.
\]
Every block in an Abelian periodic factorization of a tail of $v^\omega$ must
have the word's frequency $1/14$, so its length is a multiple of $14$. Thus
the least Abelian period length is $14$, proving sharpness in
Theorem~\ref{thm:abelian}.

\section{A frequency-denominator criterion}

We now prove Theorem~\ref{thm:denominator}. Let
\[
A(n)=\sum_{j=0}^{n-1}x_j.
\]
Having only finitely many Abelian-unbordered factors implies having only
finitely many weakly Abelian-unbordered factors, as used in the proof of
\cite[Theorem~5]{CHPZ}. Lemmas 3 and 4 of~\cite{CHPZ} therefore imply that
$x$ is balanced and has rational letter frequencies. Write the frequency of
$1$ in lowest terms as $p/q$ and define the integer discrepancy
\[
z_n=qA(n)-pn.
\]
The bounded-deviation formulation of balance in
\cite[Proposition~1]{CHPZ} implies that the set $\{z_n:n\ge0\}$ is finite.

\begin{lemma}\label{lem:border-discrepancy}
For $i\ge0$, $L\ge1$, and $1\le k\le L/2$, the factor $x[i,i+L)$ has an
Abelian border of length $k$ if and only if
\[
z_{i+k}+z_{i+L-k}=z_i+z_{i+L}.
\tag{6.1}
\]
\end{lemma}

\begin{proof}
The prefix and suffix of length $k$ contain the same number of ones exactly
when
\[
A(i+k)-A(i)=A(i+L)-A(i+L-k).
\]
Substituting $qA(n)=z_n+pn$ cancels all linear terms and gives (6.1).
\end{proof}

Put $M=\mu_{\mathrm{ab}}(x)$ and $C=M+1$. Every factor of length at least
$C$ is Abelian bordered. Choose a value $R$ which occurs infinitely often
among the $z_n$ and is maximal among the infinitely recurrent values. After
deleting a finite prefix, no value larger than $R$ occurs. Write the
successive occurrences of $R$ as
\[
s_0<s_1<s_2<\cdots
\]
and set $g_i=s_{i+1}-s_i$.

\begin{lemma}\label{lem:gaps}
For every $i$,
\[
1\le g_i<C
\qquad\text{and}\qquad
q\mid g_i.
\]
\end{lemma}

\begin{proof}
Suppose $g_i\ge C$. Apply the Abelian-border condition and
Lemma~\ref{lem:border-discrepancy} to the factor $x[s_i,s_{i+1})$. For some
$1\le k\le g_i/2$,
\[
z_{s_i+k}+z_{s_{i+1}-k}=z_{s_i}+z_{s_{i+1}}=2R.
\]
Both terms on the left are at most $R$, so both equal $R$. At least one of
the indicated positions lies strictly between $s_i$ and $s_{i+1}$,
contradicting their consecutiveness. Thus $g_i<C$.

Equal discrepancy at the endpoints gives
\[
0=z_{s_{i+1}}-z_{s_i}
=q\bigl(A(s_{i+1})-A(s_i)\bigr)-pg_i.
\]
Since $\gcd(p,q)=1$, it follows that $q\mid g_i$.
\end{proof}

\begin{proof}[Proof of Theorem~\ref{thm:denominator}]
Assume $M<2q$. Then $C=M+1\le2q$. Write $g_i=r_iq$, where $r_i$ is a
positive integer. Lemma~\ref{lem:gaps} gives
\[
0<r_i=\frac{g_i}{q}<\frac{C}{q}\le2,
\]
and hence $r_i=1$. Thus every $g_i=q$.

For every $i$, the block $x[s_i,s_{i+1})$ has length $q$, and equality of
endpoint discrepancies yields
\[
q\,|x[s_i,s_{i+1})|_1-pq=0.
\]
Every such block therefore contains exactly $p$ ones. These consecutive
blocks cover the tail beginning at $s_0$, and they all have Parikh vector
$(q-p,p)$. This is an Abelian ultimately periodic factorization with period
length $q$.
\end{proof}

When $M\ge2q$, the maximal discrepancy layer may have several return lengths
$q,2q,\ldots$. The argument above no longer selects a single block length.
Controlling only that extremal layer is not sufficient in general; the
unresolved problem requires compatibility with the other discrepancy levels.

\section{Verification and reproducibility}\label{sec:verification}

The finite results can be reconstructed from the verification material
distributed with this manuscript. The principal checks are:
\begin{enumerate}
\item exhaustive agreement of direct and prefix-sum Abelian-border tests on
all nonempty binary words of length at most $16$;
\item brute-force controls for $(C,N)=(3,6),(4,8),(5,10)$;
\item independent C++ and Python reconstruction of the $G_{14,36}$ and
$G_{15,30}$ component certificates;
\item greatest-fixed-point computation of the closed Parikh phases in the
four branching components and acyclicity of the uncovered subgraphs;
\item exact cyclic ordinary-period and Abelian-period computations, including
the two sharp witnesses.
\end{enumerate}

From the project entry directory, run
\begin{verbatim}
python -X utf8 verification\verify.py
python -X utf8 verification\verify_period_bounds.py
\end{verbatim}
The C++17 sources are \path{verification\verify_threshold14.cpp} and
\path{verification\verify_threshold15.cpp}. The programs use exhaustive
finite graph traversal; no random sampling, floating-point comparison, SAT
solver, or unverified external dataset enters the stated certificates.

\section{Scope and prior-work boundary}

The present results do not solve the full question of Charlier et al. They
establish the sharp initial ordinary-periodicity threshold, the next Abelian
periodicity layer with a sharp period bound, and a sufficient condition
depending on the denominator of the rational frequency.

A targeted search of the original paper, the 2022 survey, and indexed
follow-up literature found no result directly covering these statements.
This is an internal novelty assessment, not a claim of exhaustive global
priority. The manuscript has not undergone external expert review or formal
peer review.

Filimonova and Puzynina~\cite{FilimonovaPuzynina} recently characterized
Abelian periodicity for purely morphic binary words. Their object class and
hypotheses differ from the Abelian-border threshold conditions considered
here.

\section*{AI-assisted research and writing disclosure}

The author used OpenAI Codex for literature triage, finite-search
implementation, exact-certificate checking, proof stress-testing, and
language editing. The mathematical statements, source selection, released
artifacts, and decision to publish remain the responsibility of Carptopus.

\begin{thebibliography}{9}

\bibitem{CHPZ}
E.~Charlier, T.~Harju, S.~Puzynina and L.~Q.~Zamboni,
\newblock Abelian bordered factors and periodicity,
\newblock \emph{European J. Combin.} 51 (2016), 407--418.
\newblock \href{https://doi.org/10.1016/j.ejc.2015.07.003}{doi:10.1016/j.ejc.2015.07.003}.

\bibitem{FP}
G.~Fici and S.~Puzynina,
\newblock Abelian combinatorics on words: a survey,
\newblock \emph{Comput. Sci. Rev.} 47 (2023), 100532.
\newblock \href{https://doi.org/10.1016/j.cosrev.2022.100532}{doi:10.1016/j.cosrev.2022.100532}.

\bibitem{ES}
A.~Ehrenfeucht and D.~M.~Silberger,
\newblock Periodicity and unbordered segments of words,
\newblock \emph{Discrete Math.} 26(2) (1979), 101--109.
\newblock \href{https://doi.org/10.1016/0012-365X(79)90116-X}{doi:10.1016/0012-365X(79)90116-X}.

\bibitem{FilimonovaPuzynina}
A.~Filimonova and S.~Puzynina,
\newblock On Abelian periodicity of purely morphic words,
\newblock arXiv:2605.30306v1 (2026).
\newblock \href{https://doi.org/10.48550/arXiv.2605.30306}{doi:10.48550/arXiv.2605.30306}.

\end{thebibliography}

\end{document}
